% ============================================================================
% Fichier  : partie3/P03_C09_normes_deontologie.tex
% Titre    : Normes, Certification et Déontologie
% Auteur   : MINKA MI NGUIDJOÏ Thierry Emmanuel
% Version  : 1.0 -- novembre 2026
% Corrigés : corriges/Corr_C09.tex
% ============================================================================
% Sources : ISO/IEC 27701, 29100, 29134, 27018, 27001:2022, 27002:2022,
% ISO 19011:2018 (paraphrasées, jamais reproduites) ; loi n° 2024/017,
% art. 34, 53, 59 (texte officiel).
% Faits nouveaux (bible) : projet de certification ISO/IEC 27701 de NOUN PAY,
% septembre 2026 ; offre de KORO AUDIT à MBOLO.
% ============================================================================

\chapter{Normes, Certification et Déontologie}
\label{chap:normes}
\chaptermark{Normes et déontologie}

\epigraph{%
    « Un certificat dit qu'un système de management existe.\\
    Il ne dit pas que chaque traitement est licite. »%
}{Principe du cours}

\niveauChapitre{Fondamental}{%
    Chapitres~\ref{chap:audit} et~\ref{chap:loi}~; ISO/IEC~27001 (UE
    MSN-2S3-SEC08)
}

\begin{boxobjectifs}
\begin{itemize}
    \item Situer les normes de la protection des données~: ISO/IEC~29100,
          27701, 29134, 27018.
    \item Expliquer ce que l'ISO/IEC~27701 ajoute à l'ISO/IEC~27001.
    \item Distinguer la certification prévue par la loi (article~34) de la
          certification d'un système de management par un organisme
          accrédité.
    \item Établir une correspondance entre exigences de la loi et
          exigences d'une norme, et en dire les limites.
    \item Appliquer les règles déontologiques de l'auditeur qui manipule
          des données personnelles.
\end{itemize}
\end{boxobjectifs}

\bigskip

% ============================================================================
\section*{Avant de commencer : « avec le certificat, nous serons conformes »}
% ============================================================================

Bafoussam, mardi 1\textsuperscript{er}~septembre~2026. NOUN PAY\footnote{NOUN
PAY, ses dirigeants et ses clients sont fictifs. L'entreprise figure aussi
dans le cours de gouvernance des SI de l'auteur.} prépare une levée de fonds.
Un investisseur demande une preuve de conformité en matière de données. Le
directeur général, M.~NJOYA, a trouvé la solution~: « Nous allons nous faire
certifier ISO/IEC~27701. Avec le certificat, nous serons conformes. »

Mme~BELLO, sollicitée pour un diagnostic, lui pose deux questions. Le
périmètre du certificat couvrira-t-il la notation automatique des
crédits, ou seulement l'application de paiement~? Et l'organisme
certificateur vérifiera-t-il que chaque traitement a reçu l'autorisation de
l'article~19~? À la seconde, la réponse est non~: un auditeur de
certification vérifie le système de management contre la norme, il ne
rend pas un avis juridique.

\begin{boxalert}
\textbf{Violation typique.} Présenter un certificat de système de
management comme une preuve de conformité à la loi.

\smallskip
\textbf{Règle~:} un certificat ISO/IEC~27701 atteste qu'un système de
management de la protection des données existe et fonctionne sur un
périmètre donné. La conformité à la loi se démontre par l'audit de
conformité légale, exigence par exigence. Les deux se complètent~; l'un ne
remplace pas l'autre.
\end{boxalert}

% ============================================================================
\section{La Famille des Normes}
\label{sec:c09-famille}
% ============================================================================

\begin{tableaucours}{L{2.6cm}L{5.0cm}L{4.4cm}}{%
    \tableauHeader{Norme} &
    \tableauHeader{Objet (paraphrasé)} &
    \tableauHeader{Usage dans l'audit}
}
ISO/IEC 29100 & Cadre de la protection de la vie privée~: vocabulaire,
acteurs, onze principes & Référentiel de principes, comparaison avec la
loi \\
ISO/IEC 27701 & Système de management de la protection de la vie privée &
Audit contre une norme, certification \\
ISO/IEC 29134 & Lignes directrices pour l'évaluation d'impact & Méthode de
l'évaluation d'impact (Chapitre~\ref{chap:risque}) \\
ISO/IEC 27018 & Protection des données personnelles dans le nuage public,
pour les sous-traitants & Évaluer un hébergeur
(Chapitre~\ref{chap:transferts}) \\
ISO/IEC 27001 et 27002 & Système de management et mesures de sécurité de
l'information & Domaine sécurité (Chapitre~\ref{chap:securite}) \\
ISO 19011, ISO/IEC 27007 & Conduite des audits & Méthode
(Chapitre~\ref{chap:audit}) \\
\end{tableaucours}
\vspace{-0.8em}
\begin{center}{\small\itshape Tableau~9.1. Les normes utiles à l'audit de protection des données}\end{center}

Toutes ces normes sont publiées sous licence~: ce cours les paraphrase et
n'en reproduit aucun passage. L'auditeur qui fonde un constat sur une
exigence normative doit disposer de l'édition en vigueur et la citer avec
son année.

\subsection{Les onze principes de l'ISO/IEC 29100}

L'ISO/IEC~29100~\cite{iso29100} énonce onze principes~: le consentement et le
choix~; la légitimité et la spécification des finalités~; la limitation de
la collecte~; la minimisation des données~; la limitation de l'utilisation,
de la conservation et de la divulgation~; l'exactitude et la qualité~;
l'ouverture, la transparence et l'information~; la participation et l'accès
des personnes~; la responsabilité~; la sécurité~; et le respect des
obligations de protection de la vie privée. La comparaison avec la loi
camerounaise fait apparaître la minimisation, absente comme principe
général de la loi (Chapitre~\ref{chap:principes}).

\subsection{L'ISO/IEC 27701}

L'ISO/IEC~27701~\cite{iso27701} étend la logique de l'ISO/IEC~27001 aux
données personnelles~: contexte, leadership, planification, support,
fonctionnement, évaluation et amélioration, plus des mesures propres aux
responsables de traitement et aux sous-traitants. Publiée en 2019 comme
extension de l'ISO/IEC~27001, elle a été révisée pour devenir une norme de
système de management autonome~; l'auditeur vérifie l'édition sur laquelle
porte le certificat qu'on lui présente.

Elle apporte à l'audit trois choses que la loi ne donne pas~: une méthode
de management (un cycle d'amélioration, des revues de direction, des
audits internes), des mesures détaillées propres aux données personnelles,
et un vocabulaire partagé avec les partenaires étrangers. Elle n'apporte
pas les obligations propres au Cameroun~: l'autorisation préalable, les
trois dérogations au consentement, le rapport annuel de l'article~27, la
liste camerounaise des données sensibles.

% ============================================================================
\section{Deux Certifications}
\label{sec:c09-certification}
% ============================================================================

\begin{boxjuris}
\textbf{\loicam{n°~2024/017}, article~34}~\cite{loi2024017}~:
\textit{« L'Autorité de protection des données à caractère personnel met en
place un mécanisme de certification du traitement des données à caractère
personnel conforme aux principes et exigences énoncés par la présente
loi. »}

\smallskip
\textbf{Portée.} L'article~5 définit la certification comme un outil
permettant de communiquer sur le niveau de protection « au regard des
critères d'un référentiel approuvé au préalable par l'Autorité ». Il s'agit
d'une certification de conformité à la loi, distincte des certifications
ISO. L'article~59 punit le responsable qui « opère sans certification
prévue dans la présente loi »~: tant que le mécanisme n'existe pas, la
portée de cette sanction est un \NEK{} de critère.
\end{boxjuris}

\begin{tableaucours}{L{3.2cm}L{4.4cm}L{4.4cm}}{%
    \tableauHeader{} &
    \tableauHeader{Certification de l'article~34} &
    \tableauHeader{Certification ISO/IEC 27701}
}
Délivrée par & L'Autorité, ou selon son mécanisme & Un organisme de
certification accrédité \\
Référentiel & Approuvé par l'Autorité & Norme internationale \\
Objet & Traitement, processus, produit & Système de management sur un
périmètre \\
Portée juridique & Prévue par la loi (articles~34, 59) & Aucune au regard de
la loi \\
État en novembre~2026 & Mécanisme non publié & Disponible \\
\end{tableaucours}
\vspace{-0.8em}
\begin{center}{\small\itshape Tableau~9.2. Deux certifications à ne pas confondre}\end{center}

% ============================================================================
\section{Correspondances entre Loi et Norme}
\label{sec:c09-correspondances}
% ============================================================================

Une organisation qui veut satisfaire à la fois la loi et la norme établit une
table de correspondance. Elle évite de faire deux fois le même travail~:
le registre de l'article~29 sert aussi la norme, l'évaluation d'impact de
l'article~33 aussi, les mesures de l'article~27 aussi.

\begin{tableaucours}{L{4.4cm}L{4.0cm}L{3.6cm}}{%
    \tableauHeader{Exigence de la loi} &
    \tableauHeader{Thème de l'ISO/IEC 27701} &
    \tableauHeader{Couverture}
}
Registre (art.~29) & Registre des traitements & Complète \\
Fondement (art.~9) & Identification de la base légale & Partielle~: la
norme admet les bases locales, sans les fixer \\
Information (art.~21) & Information des personnes & Complète si les
mentions camerounaises sont ajoutées \\
Droits (art.~37 à~47) & Obligations envers les personnes & Partielle~:
délais réglementaires \\
Évaluation d'impact (art.~33) & Évaluation d'impact & Complète \\
Sécurité (art.~27) & Mesures de l'ISO/IEC 27002 & Partielle~: traçabilité
des lectures à vérifier \\
Autorisation (art.~19) & Aucun équivalent & Nulle \\
Rapport annuel (art.~27(2)) & Aucun équivalent & Nulle \\
\end{tableaucours}
\vspace{-0.8em}
\begin{center}{\small\itshape Tableau~9.3. Correspondances indicatives (thèmes paraphrasés)}\end{center}

Une correspondance n'est jamais une équivalence. Une ligne « complète » dit
que le dispositif exigé par la norme peut porter l'exigence légale~; elle ne
dit pas qu'il le fait. L'audit de conformité vérifie toujours l'exigence
légale elle-même.

% ============================================================================
\section{La Déontologie de l'Auditeur}
\label{sec:c09-deontologie}
% ============================================================================

L'auditeur de la protection des données est soumis aux principes communs de
l'audit (Chapitre~\ref{chap:audit}). Son objet ajoute trois exigences.

\subsection{Minimiser les données de l'audit}

L'auditeur ne collecte que les données nécessaires à la preuve. Il consulte
sur place plutôt que de copier, relève des indicateurs plutôt que des
contenus, masque les identités dans ses papiers de travail (« client~A,
dossier 2018-0412 »), conserve les copies indispensables dans un espace
chiffré et les détruit à la date prévue au mandat. Un auditeur qui
emporterait un extrait du fichier clients sur une clé non chiffrée
commettrait l'écart même qu'il est chargé de relever.

\subsection{L'indépendance}

L'indépendance interdit d'auditer son propre travail. KORO AUDIT a proposé
au CENTRE MÉDICAL MBOLO, en juin, de rédiger ses procédures de protection
des données~; s'il accepte, il ne pourra pas, en octobre, auditer ces
mêmes procédures sans menacer son objectivité. Mme~BELLO a refusé la
mission de rédaction et conservé la mission d'audit.

\subsection{Ce que l'auditeur découvre}

L'auditeur peut découvrir des faits graves~: un agent qui vend des fichiers,
un traitement manifestement délictueux. Il en rend compte au commanditaire,
par écrit et sans délai, comme son mandat doit le prévoir. Il ne mène pas
d'enquête qui excéderait sa mission, et il ne saisit pas lui-même la
justice, sauf obligation légale qui lui serait propre. Le
Chapitre~\ref{chap:violations} montre comment un tel fait devient une
violation à traiter, et le cours de forensique comment il devient une
enquête.

\begin{boxscenario}{Mission NOUN PAY. Le diagnostic du 8 septembre 2026}
\textbf{Constatations.} Le périmètre envisagé pour la certification couvre
l'application de paiement et exclut la notation automatique « NOUN SCORE »,
« en cours de refonte ». Aucune demande d'autorisation n'a été préparée.
L'hébergeur, établi dans l'Union européenne, présente un certificat
ISO/IEC~27018.

\textbf{Analyse} (confiance élevée). Le certificat visé laisserait hors
périmètre le traitement le plus risqué de l'entreprise, celui qui
décide des crédits. Il serait exact, et trompeur s'il était présenté comme
une preuve de conformité. Le certificat de l'hébergeur porte sur
l'hébergeur~: il ne dit rien des traitements de NOUN PAY ni du transfert
lui-même (Chapitre~\ref{chap:transferts}).
\end{boxscenario}

% ============================================================================
\section{Ce que les Normes demandent à chaque Métier}
\label{sec:c09-metiers}
% ============================================================================

\begin{boxmetier}{Délégué à la protection des données}
Bâtir le système de management sur les obligations de la loi, puis le
rapprocher de la norme~; tenir la table de correspondance et signaler les
exigences camerounaises que la norme ne porte pas.
\end{boxmetier}

\begin{boxmetier}{Auditeur de la protection des données}
Lire le périmètre et l'édition de tout certificat présenté. Ne jamais coter
\OUI{} une exigence légale sur la seule foi d'un certificat.
\end{boxmetier}

\begin{boxmetier}{Responsable de traitement}
Choisir le périmètre d'une certification en fonction des risques, pas de la
facilité~: un certificat qui exclut le traitement le plus risqué expose à un
reproche de présentation trompeuse.
\end{boxmetier}

\begin{boxresume}
\begin{itemize}
    \item L'ISO/IEC~29100 donne le vocabulaire et onze principes~;
          l'ISO/IEC~27701 un système de management~; l'ISO/IEC~29134 la
          méthode d'évaluation d'impact~; l'ISO/IEC~27018 les exigences du
          nuage.
    \item La norme n'apporte pas les obligations propres au Cameroun~:
          autorisation, dérogations au consentement, rapport annuel.
    \item La certification de l'article~34 est une certification de
          conformité à la loi, distincte d'une certification ISO.
    \item Une table de correspondance évite le double travail, mais une
          correspondance n'est pas une équivalence.
    \item L'auditeur minimise ses propres données, garde son indépendance et
          rend compte de ce qu'il découvre à son commanditaire.
\end{itemize}
\end{boxresume}

\begin{boxinfo}
\textbf{Points de contrôle.} Vous avez compris ce chapitre si vous savez~:
\begin{itemize}
    \item dire ce qu'un certificat ISO/IEC~27701 prouve et ce qu'il ne prouve
          pas~;
    \item distinguer les deux certifications~;
    \item rédiger les règles de traitement des données de l'audit.
\end{itemize}
\end{boxinfo}

\begin{boxexo}[title={\ding{45}\quad Exercice~9.1. Quelle norme \niveaubase}]
Quelle norme consultez-vous pour~: (a) conduire une évaluation d'impact~;
(b) évaluer un hébergeur en nuage~; (c) mettre en place un système de
management de la protection des données~; (d) définir le vocabulaire d'un
contrat~?
\end{boxexo}

\begin{boxexo}[title={\ding{45}\quad Exercice~9.2. Le périmètre de NOUN PAY \niveaumoyen}]
Rédigez l'avis de Mme~BELLO à M.~NJOYA sur le périmètre de la
certification, en séparant constatations et analyse. Proposez un périmètre
et un calendrier qui tiennent compte de la notation automatique.
\end{boxexo}

\begin{boxexo}[title={\ding{45}\quad Exercice~9.3. Compléter la correspondance \niveaumoyen}]
Ajoutez au Tableau~9.3 les lignes correspondant aux articles~22
(violations), 30 (sous-traitance) et 32 (transferts), avec le thème de la
norme et la couverture.
\end{boxexo}

\begin{boxexo}[title={\ding{45}\quad Exercice~9.4. Une découverte \niveauexpert}]
Pendant l'audit de MBOLO, un auditeur junior découvre qu'un agent exporte
chaque semaine la liste des patientes de la maternité vers une adresse de
messagerie personnelle. Rédigez la note au commanditaire, et dites ce que
l'auditeur doit faire et ne pas faire.
\end{boxexo}

\bigskip

\begin{boxinfo}
\textbf{Pour aller plus loin.}
\begin{itemize}
    \item L'ISO/IEC~27701~\cite{iso27701} et l'ISO/IEC~29100~\cite{iso29100}.
    \item L'ISO/IEC~27018~\cite{iso27018}, pour les sous-traitants en nuage.
\end{itemize}
\end{boxinfo}

\begin{boxinfo}
\textbf{Transition.} Le certificat de l'hébergeur de NOUN PAY pose une
question que la Partie~IV traite en premier~: comment auditer les
prestataires qui traitent des données pour le compte de l'organisation~?
C'est l'objet du Chapitre~\ref{chap:soustraitants}.
\end{boxinfo}
